\emph{Higher-dimensional non-identifiability of the diagonal moment tuple from the Gram spectrum.} Within the named constant-variance gaussian edge-block subclass of the fifth preceding paragraph, at every block dimension \(D\ge2\), write \(d_{i}=G_{ii}/\tau\) for the normalized diagonal and \((P_{2},\dots,P_{D})\) with \(P_{k}=\sum_{i}d_{i}^{k}\) for its power sums, the coordinate set carrying the \(\rho\)-family of the fourth preceding paragraph evaluated at \(d\), through the exact power-sum recursion recorded there. Pointwise, every fibre of the spectrum map over a spectrum with at least two distinct eigenvalues contains a one-parameter real-analytic curve of Gram matrices on which the diagonal moves: a plane rotation through angle \(\theta\) in the eigenbasis directions of any pair of distinct eigenvalues fixes the spectrum exactly, while along the curve \(P_{2}=\sum_{j}\lambda_{j}^{2}-2(\lambda_{i}-\lambda_{j})^{2}\sin^{2}\theta\cos^{2}\theta\) varies strictly. At \(D=3\) the diagonal image of every such fibre contains a two-dimensional patch with non-empty interior in the trace hyperplane \(\sum_{i}d_{i}=\tau\): the Jacobian of the two-angle plane-rotation construction with respect to \((d_{1},d_{2})\) is \((\lambda_{2}-\lambda_{1})(\lambda_{3}-(\lambda_{1}+\lambda_{2})/2)\) at both angles \(\pi/4\), and every non-constant three-entry spectrum admits a relabelling making this product nonzero. The rank-one spectrum \((\tau,0,\dots,0)\) itself carries the diagonal \((1,0,\dots,0)\) and the diagonal \((1/2,1/2,0,\dots,0)\) of the Gram matrix of the two-entry-equal rank-one vector, a \(D\)-independent instance of the counterexample mechanism of the first preceding paragraph. Measure-theoretically, under the named law the Gram law is invariant under unitary conjugation, so conditional on the spectrum the eigenvector matrix is Haar, and the first-row squared-moduli vector \((\lvert u_{11}\rvert^{2},\dots,\lvert u_{1D}\rvert^{2})\) is then uniform on the probability simplex - at \(D=3\) the exact simplex density \(2\) follows from the exponential change of variables \(E_{i}=T\,s_{i}\) with Jacobian \(T^{2}\) - while \(G_{11}/\tau=\sum_{j}\lambda_{j}\lvert u_{1j}\rvert^{2}\) is a linear form of that uniform vector. Every non-constant such form has a non-atomic conditional distribution: at \(D=3\) with \(\lambda_{1}>\lambda_{2}>\lambda_{3}\) and \(\mu=(\lambda_{2}-\lambda_{3})/(\lambda_{1}-\lambda_{3})\in(0,1)\), the rescaled form \(t=(G_{11}/\tau-\lambda_{3}/\tau)/(\lambda_{1}/\tau-\lambda_{3}/\tau)\) has the piecewise-linear density \(2\min(t/\mu,(1-t)/(1-\mu))\) on \((0,1)\), which integrates to \(1\), the cases \(\lambda_{1}=\lambda_{2}>\lambda_{3}\) and \(\lambda_{1}>\lambda_{2}=\lambda_{3}\) reducing to the linear \(\operatorname{Beta}(2,1)\) density \(2t\) and the mirrored density \(2(1-t)\) on \((0,1)\), while general \(D\) follows by the stick-breaking induction \((X_{1},\dots,X_{D})=(B,(1-B)X')\) with \(B\sim\operatorname{Beta}(1,D-1)\) independent of the normalized remainder \(X'\) on the \((D-1)\)-simplex, down to the \(D=2\) uniform case of the first preceding paragraph. Since the eigenvalue-collision set carries no mass under the named law, the conditional law of the diagonal tuple is non-atomic at almost every spectrum value, and because the power sums determine the diagonal only up to a permutation of \(D\) entries, the conditioning identity of the first preceding paragraph gives \(\mathbb{P}(d=F(p))=0\) for every map \(F\) of the Gram-spectrum coordinate \(p\), and the same identity with \(d\) replaced by \((P_{2},\dots,P_{D})(d)\) since a finite union of null sets is null: no map recovers the normalized diagonal tuple, its power sums, or the \(\rho\)-family evaluated at the diagonal from the Gram-spectrum coordinate almost everywhere on the support of the law of any named constant-variance edge-block datum at any \(D\ge2\), extending the \(D=2\) single-statistic nonexistence of the first preceding paragraph to the full diagonal moment tuple at every dimension; the identification of the fibre diagonal image with the full Schur--Horn permutahedron is outside the proved content of this paragraph. These statements are exact linear-algebraic and measure-theoretic identifiability statements on the support of the named subclass at every \(D\ge2\), with the fibre-patch content at \(D=3\), they bound nonexistence between the named coordinates only, they evaluate no matrix integral of the core action, they supply no state and no selection mechanism, they do not fix the state that sets \(W\), \(g\), \(s\) or the frame, they leave the missing steps (i)-(iv) of the Missing step paragraph unchanged, and they close nothing in Problem~G.
